\documentclass{article}
% Shared layout. Keep benchmark prose and result-specific plots outside this folder.
\PassOptionsToPackage{numbers}{natbib}
\usepackage[preprint]{neurips_2026}
\usepackage[T1]{fontenc}
\usepackage[hidelinks]{hyperref}
\usepackage{booktabs,amsmath,microtype,pgfplots}
\pgfplotsset{compat=1.17}
\usetikzlibrary{patterns}
\definecolor{okblue}{HTML}{0072B2}
\definecolor{okorange}{HTML}{D55E00}
\pgfplotsset{every axis/.append style={font=\small, tick label style={font=\small}, width=\linewidth, height=5cm, grid=major}}

\newcommand{\PaperTitle}{Benchmark Ranking Stability: Separating Score Order from Evidence}
\newcommand{\PaperAbstract}{A benchmark score orders systems without showing how that order depends on the evaluation sample. Benchmark Ranking Stability combines paired item resampling, repeated runs, and prompt-specific ranks. Its public suite contains 60 reasoning items across 6 families, each with 5 prompts. We run 9 scripted or mock controls twice through the actual evaluator. We also evaluate simulated populations with different sizes and correctness gaps, using 12 seeds per condition. The answer-key control reaches 1.000 [0.940, 1.000] all-variant item success (Wilson interval). In simulated populations of eight agents, a designated leader's mean bootstrap first-place share changes from 0.150 [0.090, 0.222] at equal ability to 0.644 [0.568, 0.728] at the largest gap. These are confidence intervals across simulation seeds. Exact ties retain high rank agreement while offering no unique winner. The contribution is an inspectable ranking diagnostic and controlled validation, not a learned-model leaderboard or a new estimator.}
\newcommand{\AuthorOrcid}{0009-0005-9964-6841}
\newcommand{\RivalProb}{0.550}
\newcommand{\Items}{60}
\newcommand{\Tasks}{300}
\newcommand{\Variants}{5}
\newcommand{\Families}{6}
\newcommand{\Seeds}{12}
\newcommand{\Runs}{3}
\newcommand{\Draws}{500}
\newcommand{\Level}{95}
\newcommand{\SimReports}{1512}
\newcommand{\PolicyCount}{9}
\newcommand{\PolicyExecutions}{18}
\newcommand{\ApiCost}{0}
\newcommand{\AblationShare}{0.051 [0.020, 0.080]}
\newcommand{\AblationWidth}{-0.167 [-0.417, 0.000]}
\newcommand{\TieTau}{1.000}
\newcommand{\TieRank}{1.500}
\newcommand{\TieShare}{0.500}

\title{\PaperTitle}
\author{Pooria Arab\\Independent researcher\\ORCID: \AuthorOrcid}
\begin{document}
\maketitle
\begin{abstract}\PaperAbstract\end{abstract}
\section{Why ranking stability matters}
A model can lead an evaluation without having a reliable lead. Its position depends on which questions appear, which prompt introduces them, and which run supplies the answers. A leaderboard often hides these choices behind a single score. That score describes the observed sample. It does not by itself show whether nearby alternatives would change the winner. This distinction matters when a team selects a system, compares releases, or decides which result deserves further study.\par
Benchmark Ranking Stability makes this dependence the object of evaluation. Its construct is \emph{conditional ranking stability}: the persistence of score order when the supplied task sample and recorded runs are resampled. A separate prompt analysis shows whether the ordering changes across the supplied prompt forms. The construct is narrower than general reasoning ability. It also differs from answer consistency, which can remain high for a system that gives the same wrong answer repeatedly.\par
The benchmark couples a ranking analysis with a small reasoning suite. The analysis can compare external reports that share task identifiers. The suite supplies controlled items with several prompt forms. These components address different questions. The suite asks whether an answer meets a deterministic key. The analysis asks how much evidence supports the ordering of the resulting scores. Neither component makes a deployment decision on its own.\par
Our contribution is a controlled validation of this combination. Scripted policies establish a floor and a key-assisted ceiling. A simulated population sweep varies the number of rivals and the leader's correctness advantage. A simulated analysis ablation removes item grouping while preserving all answers. Exact-tie and simulated prompt controls distinguish stable order from an informative winner. These experiments use actual evaluation and analysis paths, but they do not measure trained language models.\par
A useful reading rule follows: report the score, the ranking uncertainty, and the comparison set together. A first-place share means little without the set of competitors. A narrow rank interval can conceal a broad uncertainty in the score gap. High rank agreement can also arise from permanent ties. We therefore present the metrics as complementary views of a conditional experiment, rather than interchangeable measures of model quality.\par
\section{Related work}
The statistical foundations are established. Efron introduces bootstrap resampling as a way to study uncertainty from observed data \citep{efron}. Kendall develops a treatment of ties for rank correlation \citep{kendall}. This benchmark applies those ideas to reports that include tasks, prompt variants, and repeated runs. It does not propose a new estimator or establish a new coverage theorem. The empirical question is whether the resulting diagnostic behaves usefully under explicit controls.\par
Miller treats language-model evaluations as experiments and emphasizes uncertainty, paired comparisons, and sampling assumptions \citep{miller}. Our study follows that framing by stating what each interval varies. Task uncertainty, run variation, and simulation replication are separate levels. We do not combine repeated executions of a deterministic policy into a larger task sample. Nor do we interpret repeated bootstrap draws as additional measured model responses.\par
Prompt sensitivity provides another motivation. Sclar and colleagues study performance changes under meaning-preserving prompt formatting and propose evaluating a range of formats \citep{sclar}. Our prompt controls have a different purpose. They deliberately encode a dependence on prompt identity so that a ranking reversal has a known cause. This is a check on the measurement surface, not evidence that real systems respond to these prompts in the same way.\par
The family study examines how benchmark construction and probing reveal evaluation defects \citep{family}. Here the unit of study is the ranking benchmark itself. We ask what its reported statistics mean and where their interpretation fails. The benchmark's reasoning bank supports controlled experiments; it is not a representative collection of human requests. Comparisons with broad capability suites therefore need a scope qualification. Broad coverage and explicit ranking uncertainty solve different problems, and both can be necessary.\par
\section{Experiment design}
\subsection{Observation units and matched conditions}
The public bank contains \Items{} underlying items and \Tasks{} prompt rows. Each item has \Variants{} prompt forms. All forms preserve the central problem statement and expected answer while changing the instruction around it. This creates matched observations: differences across forms can be inspected without replacing the underlying question. These are matched wording conditions, not adversarial traps paired with safe lookalikes.\par
An observation is a graded answer for one item, prompt form, agent, and run. An item cluster contains all forms of that item. The cluster is the sampling unit for the ranking bootstrap. The complete cluster travels together when it is selected. The same cluster draw applies to every compared agent, which preserves task pairing. A separate random selection chooses one observed run for each agent in each resample.\par
The recorded means use every supplied run. The resampled ranking instead uses one selected run per agent. Thus its variability describes task resampling and a recorded-run choice. It is not the sampling distribution of a newly averaged set of runs. This distinction becomes important when a future report averages many independent calls before ranking. Such a report targets a different decision rule and needs a corresponding resampling design.\par
\subsection{Scripted controls}
We evaluate \PolicyCount{} controls, each with an independent execution repeat. Empty and refusal policies produce no substantive answer. Numeric constants test incidental credit. Echo returns the prompt unchanged. First-number and last-number policies extract a visible number without solving the problem. The answer-key policy returns the expected answer and defines a privileged ceiling. The packaged mock supplies a deterministic mixture of correct and incorrect responses.\par
Weak policies receive prompt text only. The answer key is available to the ceiling policy and to the simulated answer generator by design. This separation prevents the ceiling from being described as an ordinary solver. The mock also serves as an execution control, not a trained agent. All controls use the public evaluator. We retain the answers, scores, and reasons, then check the semantic reports across repeated executions.\par
The floor question is deliberately modest: do these named weak policies obtain substantial credit on this fixed bank? A low score cannot prove that every shortcut fails. Echo can receive credit when the prompt contains an answer-like value or option mention. A constant can match a key by accident. Those outcomes remain in the result. We do not remove items after seeing which policy succeeds.\par
\subsection{Simulated population sweep}
The additional virtual experiment is explicitly simulated. Each synthetic agent has a fixed correctness probability. Rivals share correctness probability $\RivalProb$; a designated leader has an added gap. The sweep varies population size and that gap, as shown in the result table. Every cell uses \Seeds{} simulation seeds and \Runs{} recorded runs per agent. Within each run, one Bernoulli draw determines correctness for an item, and all its prompt forms receive the same outcome.\par
Correct simulated answers use the expected key. Incorrect simulated answers are empty. Both pass through the actual evaluator. This construction isolates ranking behavior from language generation and answer parsing. It assumes complete within-item correlation, independent errors across agents, and independent errors across recorded runs. It does not include family-specific difficulty, shared training data, correlated model failures, or systematic reasoning strategies.\par
The random sequence is held fixed across gap conditions within a population and seed. Raising the leader's threshold therefore preserves already-correct leader items and can add new correct items. The rivals' answers remain fixed. This pairing supports a direct comparison of gaps without adding unrelated simulation noise. The leader is named before sampling, including in the equal-ability condition. We never select the empirical winner and then label it the leader.\par
\subsection{Control-layer ablation and diagnostic cases}
The simulated ablation uses the same reports from the smallest population at the middle gap. The normal analysis groups prompt forms into items. The ablated analysis removes the grouping information, so each prompt row becomes a sampling unit. It preserves answers, scores, agents, and recorded runs. The intervention therefore tests a statistical control layer: recognition that repeated forms of an item are correlated.\par
A separate simulated prompt control creates opposing preferences. One agent succeeds on the early prompt forms; the other succeeds on the remaining forms. The expected per-form rank reversal follows directly from that construction. An exact-tie control gives distinct agent labels to identical key-assisted reports. Together these controls check that the analysis can display both prompt dependence and permanent equality without inventing a unique winner.\par
\section{Coverage and task scope}
\begin{center}\small\begin{tabular}{lll}\toprule
Family & Items & Prompt rows \\\midrule
multi-step-arithmetic & 10 & 50 \\
sequence & 10 & 50 \\
remainder & 10 & 50 \\
ordering & 10 & 50 \\
set-counting & 10 & 50 \\
rule-tracing & 10 & 50 \\
\bottomrule\end{tabular}\end{center}

The bank spans \Families{} compact reasoning families. Arithmetic items require a sequence of operations. Sequence items extend a numerical pattern. Remainder items require modular arithmetic. Set-counting items combine disjoint and overlapping groups according to their wording. Rule-tracing items apply conditional transformations. Ordering items identify an option from before-and-after relations. The balanced family counts are a design property, not a measurement of real usage frequency.\par
These families offer short, checkable answers and make repeat experiments inexpensive. They also impose strong limits. The benchmark contains no open-ended writing, retrieval corpus, tool negotiation, long conversation, or changing external environment. It cannot establish how a rank transfers to those settings. A stable ordering here can coexist with an unstable ordering elsewhere because the underlying task distribution changes.\par
The prompt forms vary the instruction surrounding a fixed problem. They do not exhaust the space of valid paraphrases. They also do not change all aspects of prompt presentation independently. A measured difference cannot automatically be assigned to politeness, length, reasoning cues, or answer-format demands. Such attribution requires a separate design that changes one factor while holding others constant.\par
The public bank is visible and can be memorized. This study uses it to validate measurement behavior, so visibility is acceptable for the stated purpose. A learned-model comparison would need stronger protection against contamination and adaptive prompt selection. Fresh tasks can reduce direct recall but do not automatically provide representativeness. The intended use population still needs a definition and a justified sampling process.\par
\section{Taxonomy of ranking failures}
\paragraph{Task-sample reversal.} Two agents can succeed on different subsets while having similar overall means. Changing which items receive weight can reverse their order. This is not necessarily a model failure. It is evidence that the observed lead depends on the sample. A paired resample preserves the fact that both agents faced the same selected items.\par
\paragraph{Run-choice reversal.} An agent can vary between recorded executions even when the tasks stay fixed. Selecting a different run may change the ranking. The current analysis represents this with a single run draw per agent. It cannot estimate unobserved changes in provider behavior, decoding settings, or future releases. If those conditions differ between reports, the comparison mixes effects that the analysis cannot separate.\par
\paragraph{Prompt-conditioned reversal.} Two agents can prefer different instruction forms. Averaging those forms creates an overall ordering that may hold for none of the intended uses. The prompt ranking table exposes the supplied forms individually. Its best-to-worst rank range is descriptive. It is not a confidence interval over the unbounded population of future prompts.\par
\paragraph{False precision from duplicated evidence.} Treating correlated prompt forms as independent observations increases the apparent sample size. The answers do not become more diverse merely because they occupy separate rows. Item grouping protects against this error when the identifiers correctly encode the grouping. If related items lack that structure, the analysis cannot infer the missing dependence automatically.\par
\paragraph{Stable but uninformative ties.} Equal scores can remain equal under every resample. Rank agreement then appears perfect, although no unique winner exists. Splitting first-place credit and reporting average tied ranks prevents an alphabetical ordering from masquerading as evidence. The exact-tie control illustrates why agreement should be read beside first-place share and the underlying score distribution.\par
\paragraph{Selection and alignment loss.} Adding rivals changes what it means to rank first. Dropping missing tasks changes the evaluated task population. Neither change requires a modification to the agents themselves. Reports must therefore state the comparison set and the shared item count. A filtered intersection can favor one agent if missingness depends on difficulty or failure.\par
\section{Metrics and scoring}
Let $y_{mirv}$ denote the score for agent $m$, item $i$, run $r$, and prompt form $v$. The complete bank yields binary scores after deterministic answer normalization. For balanced reports, the mean is $\bar y_m=(IRV)^{-1}\sum_{i,r,v} y_{mirv}$. External reports can supply scores between zero and one. Ranking uses decreasing means, with average ranks for exact ties. A ranking is therefore relative to the supplied agents and shared tasks.\par
The grader strips a surrounding code fence and can read a final JSON answer object. It normalizes surface form before comparison. Numeric answers use the first number token. Ordering answers use option mentions and order cues. These rules accept several answer formats, but they are not a semantic proof. A correct value mentioned for the wrong reason can pass, and an explanation with an earlier incidental number can fail.\par
For each bootstrap draw, item indices are sampled with replacement. All prompt forms of a selected item remain together. One recorded run is selected for each agent, and the sampled means determine new ranks. We use \Draws{} draws per analysis. The reported rank interval is the central \Level{} percent of these ranks. It describes this resampling procedure; it does not certify frequentist coverage for an unknown universal model order.\par
First-place share assigns equal credit to agents tied at the best sampled score. If $W_b$ is the set of winners in draw $b$, the share for agent $m$ is $B^{-1}\sum_b \mathbf{1}[m\in W_b]/|W_b|$. Fractional ties make a simple binomial interval inappropriate. In the simulated sweep, uncertainty bars instead resample the independent simulation seeds and summarize the mean share across those seeds.\par
Rank agreement uses Kendall's tau-b. With concordant pairs $C$, discordant pairs $D$, and pairs tied on only one side $T_x,T_y$, the coefficient is $(C-D)/\sqrt{(C+D+T_x)(C+D+T_y)}$. The analysis averages agreement between the full-data ranking and each sampled ranking. When both rank vectors are identical and wholly tied, the implementation returns perfect agreement by convention. That special case is disclosed because the usual denominator vanishes.\par
All confidence intervals use nominal \Level{} percent bounds. Policy mean intervals resample complete item clusters. They are conditional descriptions under an exchangeable-item interpretation of this fixed bank. They are not uncertainty from repeated model calls. At boundary policies, the empirical bootstrap cannot produce unseen successes or failures, so its interval collapses. We additionally report Wilson intervals for the binary event that every prompt form of an item is correct. This is a different estimand and uses items as its denominator.\par
For simulated cells, confidence intervals resample the \Seeds{} independent seed summaries. They describe uncertainty in a simulation replication mean. Paired gap and ablation intervals resample paired seed differences. The seed count is small, and percentile intervals may be rough. No formal hypothesis test or multiplicity-adjusted claim is made. We report the full planned grid and treat the intervals as descriptive evidence.\par
\section{Results}
\subsection{Scripted floor, privileged ceiling, and repeatability}
\begin{center}\small\begin{tabular}{lll}\toprule
Control & Mean score (cluster CI) & All variants correct (Wilson CI) \\\midrule
empty & 0.000 [0.000, 0.000] & 0.000 [0.000, 0.060] \\
zero & 0.000 [0.000, 0.000] & 0.000 [0.000, 0.060] \\
forty-two & 0.033 [0.000, 0.083] & 0.033 [0.009, 0.114] \\
echo & 0.033 [0.000, 0.083] & 0.033 [0.009, 0.114] \\
refusal & 0.000 [0.000, 0.000] & 0.000 [0.000, 0.060] \\
first-number & 0.000 [0.000, 0.000] & 0.000 [0.000, 0.060] \\
last-number & 0.000 [0.000, 0.000] & 0.000 [0.000, 0.060] \\
answer-key & 1.000 [1.000, 1.000] & 1.000 [0.940, 1.000] \\
mock & 0.500 [0.473, 0.527] & 0.000 [0.000, 0.060] \\
\bottomrule\end{tabular}\end{center}

\begin{figure}[ht]\centering\begin{tikzpicture}\begin{axis}[ymin=-0.05,ymax=1.1,ylabel={Mean score},xtick={0,1,2,3,4,5,6,7,8},xticklabels={empty,zero,42,echo,refusal,first,last,key,mock},x tick label style={rotate=35,anchor=east}]
\addplot+[only marks,mark=*,color=okblue,error bars/.cd,y dir=both,y explicit] table[x=index,y=score,y error minus=low,y error plus=high,col sep=comma]{data/policies.generated.csv};
\end{axis}\end{tikzpicture}\caption{Scripted and mock controls. Bars show item-cluster bootstrap confidence intervals. The key control has access to expected answers. Boundary intervals collapse; the accompanying table reports Wilson intervals for all-variant item success. First and last refer to number extraction.}\label{fig:policies}\end{figure}

The key-assisted policy reaches the ceiling, while every named weak policy remains far below it. This supports the planned floor-versus-ceiling check. The nonzero constant and echo scores also show that incidental credit exists. They remain visible in the table rather than being explained away as invalid submissions. The result applies to this fixed policy set and public bank; it is not a universal shortcut-resistance guarantee.\par
The mock's item-level all-variant outcome differs sharply from its mean answer score. A policy can obtain substantial credit across rows while failing to solve every prompt form of any item. This difference demonstrates why a repeated-form reliability measure cannot be replaced by a row mean. Conversely, the Wilson measure deliberately discards partial within-item success and should not replace the primary score without an explicit purpose.\par
All \PolicyExecutions{} primary control executions complete without a reported evaluator error. Each policy and its repeat have identical semantic traces. Timing and start timestamps are excluded from that comparison. This confirms repeatability for the chosen deterministic controls. It does not add independent items, estimate stochastic model reliability, or validate the grader against human judgments. The recorded API cost is \ApiCost{} dollars.\par
\subsection{Simulated population and gap effects}
\begin{center}\small\begin{tabular}{llll}\toprule
Agents & Gap & Leader share (CI) & Tau-b (CI) \\\midrule
2 & 0.000 & 0.574 [0.459, 0.693] & 0.391 [0.268, 0.507] \\
2 & 0.050 & 0.725 [0.629, 0.822] & 0.486 [0.338, 0.651] \\
2 & 0.150 & 0.922 [0.878, 0.961] & 0.844 [0.756, 0.922] \\
4 & 0.000 & 0.313 [0.211, 0.424] & 0.345 [0.279, 0.413] \\
4 & 0.050 & 0.475 [0.364, 0.596] & 0.380 [0.306, 0.456] \\
4 & 0.150 & 0.791 [0.726, 0.857] & 0.570 [0.533, 0.608] \\
8 & 0.000 & 0.150 [0.090, 0.222] & 0.291 [0.245, 0.348] \\
8 & 0.050 & 0.289 [0.207, 0.384] & 0.312 [0.260, 0.373] \\
8 & 0.150 & 0.644 [0.568, 0.728] & 0.413 [0.379, 0.458] \\
\bottomrule\end{tabular}\end{center}

\begin{figure}[ht]\centering\begin{tikzpicture}\begin{axis}[xlabel={Simulated leader probability gap},ylabel={Mean bootstrap first-place share},ymin=0,ymax=1.1,legend pos=north west]
\addplot+[mark=*,error bars/.cd,y dir=both,y explicit] table[x=gap,y=share,y error minus=low,y error plus=high,col sep=comma]{data/population-2.generated.csv};\addlegendentry{Two simulated agents}
\addplot+[mark=square*,error bars/.cd,y dir=both,y explicit] table[x=gap,y=share,y error minus=low,y error plus=high,col sep=comma]{data/population-4.generated.csv};\addlegendentry{Four simulated agents}
\addplot+[mark=triangle*,error bars/.cd,y dir=both,y explicit] table[x=gap,y=share,y error minus=low,y error plus=high,col sep=comma]{data/population-8.generated.csv};\addlegendentry{Eight simulated agents}
\end{axis}\end{tikzpicture}\caption{Simulated population sweep. Bars are confidence intervals for the mean across independent simulation seeds. They are not confidence intervals for real agents. The designated leader is fixed before generating results; it is not selected from observed winners.}\label{fig:populations}\end{figure}

The simulated designated leader's mean first-place share rises as its correctness advantage grows. At each gap, the share falls as more rivals enter the comparison. These are the planned directional patterns. The equal-ability cells do not produce identical sample shares for every agent. Finite simulations can favor the designated agent by chance, and the intervals show the resulting uncertainty. Equal generating probabilities do not force identical finite reports.\par
\begin{center}\small\begin{tabular}{ll}\toprule
Simulated agents & Paired share difference (CI) \\\midrule
2 & 0.348 [0.259, 0.436] \\
4 & 0.479 [0.415, 0.551] \\
8 & 0.494 [0.449, 0.552] \\
\bottomrule\end{tabular}\end{center}

The paired contrasts compare the largest gap with equal ability within the same simulated population. Their intervals preserve the shared random draws across conditions. They describe the response of this simulation mechanism to an imposed advantage. They do not estimate the benefit of training a real model or changing its prompt. No real intervention creates the gap in this study; the probability threshold defines it.\par
The simulated tau-b results answer a different question from the leader share. With more rivals, many comparisons occur among equal-ability agents. Their order can vary even when the leader becomes more distinct. A stronger leader therefore need not produce near-perfect agreement for the full ranking. This is a reason to choose a metric according to the decision: selecting a winner and ordering all alternatives are different tasks.\par
The simulated grid contains \SimReports{} evaluated reports. Multiple runs allow run-choice variation to enter the bootstrap. The whole grid remains synthetic, despite using actual scoring and analysis code. Its evidence concerns the behavior of the diagnostic under a known response process. It cannot establish that the chosen probability levels, independence assumptions, or within-item correlation match deployed systems.\par
\subsection{Simulated grouping ablation}
\begin{figure}[ht]\centering\begin{tikzpicture}\begin{axis}[xlabel={Simulation seed},ylabel={Leader rank interval width},ymin=-0.05,ymax=1.2,legend pos=south west]
\addplot+[mark=*,color=okblue] table[x=seed,y=clustered,col sep=comma]{data/ablation.generated.csv};\addlegendentry{Item clusters}
\addplot+[mark=square*,color=okorange,dashed] table[x=seed,y=independent,col sep=comma]{data/ablation.generated.csv};\addlegendentry{Prompt rows treated independently}
\end{axis}\end{tikzpicture}\caption{Simulated control-layer ablation on paired reports. Each point is a conditional rank interval width, not a success rate. Removing item grouping can change the analysis without changing any answer. Overlapping lines indicate a null at the rank-width resolution.}\label{fig:ablation}\end{figure}

In the simulated ablation, removing item grouping changes the leader's mean first-place share by \AblationShare{}. The interval summarizes paired independent-minus-clustered differences across seeds. The answers and their scores remain the same. The observed shift therefore comes from how the analysis samples evidence, rather than improved agent performance. It supports retaining item grouping when repeated forms share an underlying response.\par
The simulated rank-width difference is \AblationWidth{}. Its interval includes no change. This is a null at the chosen summary resolution, and it remains a null in our interpretation. Rank intervals are discrete; for a small population, many different sampled distributions yield the same endpoints. A first-place share can move while the rank interval stays unchanged. The ablation does not justify a claim that every reported uncertainty measure detects the grouping error.\par
The ablated analysis still samples recorded runs, which preserves another source of uncertainty. It is therefore not a clean comparison between complete dependence and complete independence of all observations. It removes one control layer while retaining the rest. A broader ablation would vary run aggregation, item structure, and prompt dependence separately. This study does not infer those unmeasured effects from the present contrast.\par
\subsection{Exact ties and simulated prompt reversals}
The exact-tie control yields mean tau-b \TieTau{} and average rank \TieRank{} for each agent. Both agents split first place equally, with share \TieShare{} and conditional interval $[\TieShare{},\TieShare{}]$. This equality is deterministic: its conditional uncertainty range is the same point. The result is a warning about interpretation, not a statistical discovery. Perfect agreement here means persistent equality, not evidence for either agent as a unique winner.\par
\begin{center}\small\begin{tabular}{lll}\toprule
Prompt variant & prompt-1 & prompt-0 \\\midrule
1 & 2 & 1 \\
2 & 2 & 1 \\
3 & 1 & 2 \\
4 & 1 & 2 \\
5 & 1 & 2 \\
\bottomrule\end{tabular}\end{center}

The simulated prompt control produces the constructed reversal across prompt forms. The table reports ranks rather than artificial accuracy estimates. Both agents have complete item coverage under every form, so the reversal cannot be attributed to missing tasks. The pooled ranking weights each included form equally. Changing the number or weight of forms could change the pooled result without changing behavior on any existing form.\par
These controls show why no single statistic suffices. The simulated population sweep has uncertainty from sampled successes and recorded runs. The exact-tie control has no distinct winner. The simulated prompt control has form-specific preferences. A reporting format that shows only the best mean score would hide all three differences. The combined diagnostic makes those distinctions inspectable without claiming to resolve their practical importance.\par
\section{Roadmap}
The next empirical step is a learned-model study with a fixed comparison roster, unchanged provider settings, and repeated calls. It should preregister task sampling, decoding, exclusions, and a stopping rule before any scored response. A free inference route can support that study only if its usage and publication terms permit the work. No learned-model performance is projected from the simulated cells reported here.\par
A second priority is validation of the answer interpretation rules. Independent reviewers should examine accepted wrong answers and rejected correct answers across families and output forms. That work should separate task-key errors from parsing errors and reasoning failures. The present key-assisted ceiling cannot detect an incorrect key, because it uses that key as the answer. Human validation remains a distinct missing source of evidence.\par
A statistical extension should test interval coverage under several known response processes. Useful conditions include correlated agents, unequal item difficulty, partial prompt correlation, and informative missingness. Coverage needs repeated synthetic datasets with a defined population target. Re-running the same deterministic trace does not supply that evidence. The current simulated sweep demonstrates sensitivity and controlled contrasts, not calibrated coverage across all such mechanisms.\par
Future prompt coverage should also separate instruction wording from formatting and answer constraints. Independently sampled paraphrases would allow more direct generalization beyond the fixed forms. A study should state whether it targets a preferred prompt, an average prompt, or worst-case performance over a defined set. Those targets require different summaries and should not be selected after inspecting the results.\par
\section{Caveats and intended use}
This is a controlled measurement study, not a leaderboard. We run no trained language model and collect no human survey or human performance baseline. Free hosted inference is not configured for this study, and we do not probe provider availability. No model-provider scores are published. The absence of that experiment limits external validity; it is not evidence that free models are unavailable elsewhere.\par
The local preregistration is committed before result collection, but it is not an independently timestamped public registration. The study also uses a public bank whose design and expected answers are known. These choices suit a reproducible diagnostic study but provide weaker protection against researcher flexibility than an external registration and held-out evaluation. The fixed grid and explicit null reporting reduce that flexibility without eliminating it.\par
The item bootstrap treats observed items as exchangeable. The bank is balanced by family, while the bootstrap can change family proportions. If a target population requires a fixed family mix, stratified sampling would answer a different and potentially more useful question. The present analysis does not estimate how family weights should be chosen. Those weights are a substantive decision about intended use.\par
Exact equality uses finite numeric scores. Near ties are not automatically treated as practical equivalence. A tiny mean difference can receive a strict ordering even when its practical effect is negligible. Conversely, a tied score can hide different error patterns. Users should inspect paired outcomes and a meaningful score-gap threshold before interpreting a rank as a reason to select one system over another.\par
The shared-task intersection is another boundary. A report with missing hard items can change which evidence remains for every agent. The analysis checks that usable shared data exist, but it cannot reconstruct missing outcomes or justify their removal. This study has complete coverage and does not experimentally quantify missingness bias. Its conclusions should not be transferred to incomplete external reports without a separate audit.\par
Simulation gives precise control over the source of variation at the cost of realism. All prompt forms share a simulated correctness event within an item. Real errors may be partly shared, systematically related across families, or correlated across agents. A simulated first-place share is therefore conditional on the stated mechanism. It is not a probability that a real model is the best system for a user.\par
\section{Citation, availability, and authors}
Code, generated results, and reproduction instructions are available at \url{https://github.com/pooriaarab/benchmark-ranking-stability}. The paper and generated figures use CC-BY-4.0. The benchmark retains its existing software and task license. Cite this paper by its title and author; no new venue acceptance or publication DOI is claimed.\par
Pooria Arab is an independent researcher. The study uses automated assistance for implementation and drafting. The author remains responsible for the scientific claims and release decisions. The evidence is limited to the scripted controls, the labelled simulations, and the repeat checks described here. No participant recruitment or private user data enters the experiments.\par
\bibliographystyle{plainnat}
\bibliography{references}
\end{document}
